\documentclass[11pt]{article}

\usepackage[utf8]{inputenc}
\usepackage[T1]{fontenc}
\usepackage{lmodern}
\usepackage[margin=1in]{geometry}
\usepackage{amsmath, amssymb, amsfonts, mathtools}
\usepackage{booktabs}
\usepackage{array}
\usepackage{enumitem}
\usepackage{hyperref}

\hypersetup{
  colorlinks=true,
  linkcolor=blue,
  urlcolor=blue,
  pdftitle={KK and FR implementations in the grip package},
  pdfauthor={Codex}
}

\setlength{\parskip}{0.6em}
\setlength{\parindent}{0pt}

\newcommand{\R}{\mathbb{R}}
\newcommand{\vect}[1]{\boldsymbol{#1}}
\newcommand{\norm}[1]{\left\lVert #1 \right\rVert}

\title{How KK- and FR-style Layouts Are Implemented in the \texttt{grip} Package\\
\large A Code-Based Comparison with the Original Kamada--Kawai and Fruchterman--Reingold Algorithms}
\author{Generated from direct inspection of the current \texttt{grip} source tree}
\date{\today}

\begin{document}

\maketitle

\begin{abstract}
This note describes precisely how KK- and FR-style layout updates are implemented
inside the current \texttt{grip} package. The short answer to the motivating
question is: yes, the package does \emph{not} implement the textbook full
all-pairs Kamada--Kawai algorithm, and it also does \emph{not} implement the
textbook full all-pairs Fruchterman--Reingold algorithm. Instead, the package
uses a multiscale GRIP engine whose active kernels are local KK-style and
local FR-style direction computations, optionally augmented by an exact-or-sampled
active-set repulsion term. This document explains the exact code paths, the
mathematical formulas implied by the code, and the differences from the original
KK and FR methods.
\end{abstract}

\section{Executive summary}

The most important conclusions are:
\begin{enumerate}[leftmargin=1.5em]
\item \textbf{There is no full original KK implementation in the package core.}
The code does not build an all-pairs shortest-path matrix, does not define the
classical KK spring energy over all vertex pairs, and does not perform the
classical Newton--Raphson per-vertex solve.

\item \textbf{There is no full original FR implementation in the package core.}
The active FR kernel uses edge-based attraction, but repulsion is restricted to
a retained local BFS neighborhood rather than all other vertices.

\item \textbf{The package has several KK-style routines, not just one.} The
active code contains:
  \begin{itemize}[leftmargin=1.5em]
  \item \texttt{KK\_spring\_local()}, used only for warm-starting newly introduced vertices,
  \item \texttt{KK\_spring\_v4()}, used on coarse and intermediate MISF levels,
  \item \texttt{KK\_spring\_final()}, used on the finest level only when
  \texttt{final\_mode = "kk\_repulse"}.
  \end{itemize}

\item \textbf{The extra long-range interaction in current \texttt{grip.layout()}
is repulsion only.} It is not a full all-pairs KK spring model. The added term
is an active-set repulsion over currently active vertices, computed either
exactly or by uniform sampling with a scale correction.

\item \textbf{The only FR routine used by the engine is \texttt{FR\_spring()}.}
There is also a function \texttt{FR\_spring\_v2()}, but it is not called by the
current engine.
\end{enumerate}

\section{Code paths and public entry points}

The public R-level entry points map to two main C++ wrappers:
\begin{itemize}[leftmargin=1.5em]
\item \texttt{grip.layout.legacy()} calls \texttt{grip\_layout\_adj\_cpp()}.
\item \texttt{grip.layout()} forwards to \texttt{grip.layout.globalrep()}.
\item \texttt{grip.layout.globalrep()} calls \texttt{grip\_layout\_globalrep\_adj\_cpp()}.
\end{itemize}

The relevant source paths are:
\begin{itemize}[leftmargin=1.5em]
\item \texttt{R/grip\_layout.R}
\item \texttt{src/grip\_rcpp.cpp}
\item \texttt{src/DrawGraph.h}
\item \texttt{src/DrawGraph.cpp}
\item \texttt{src/MishEngine.cpp}
\item \texttt{src/MishSupport.cpp}
\end{itemize}

\subsection{What the R wrappers expose}

At the R level, the current primary API is:
\[
\texttt{grip.layout(..., final\_mode = c("fr", "kk\_repulse"),}
\]
\[
\texttt{coarse\_repulsion\_factor, coarse\_repulsion\_sample,}
\]
\[
\texttt{coarse\_repulsion\_exact\_below, final\_anchor\_factor,}
\]
\[
\texttt{final\_move\_scale\_after\_first).}
\]

From the source:
\begin{itemize}[leftmargin=1.5em]
\item \texttt{grip.layout()} is just a forwarder to
\texttt{grip.layout.globalrep()}.
\item \texttt{grip.layout.legacy()} preserves the historical local-force behavior.
\item \texttt{final\_mode = "fr"} keeps the FR-style finest-level stage.
\item \texttt{final\_mode = "kk\_repulse"} replaces that finest-level FR stage by
a KK-style local update plus active-set repulsion.
\end{itemize}

\section{The shared GRIP engine structure}

\subsection{MISF hierarchy}

The engine is not ``run KK'' or ``run FR'' in the textbook sense. It is a
multiscale GRIP engine built on a maximal independent set filtration (MISF).
Vertices are introduced progressively from coarser to finer active sets.

At each active level:
\begin{enumerate}[leftmargin=1.5em]
\item a subset of vertices is currently active,
\item each active vertex has a retained neighborhood list,
\item the engine applies a local force-style update routine to each active vertex,
\item the resulting direction is converted into a move whose magnitude is controlled by
a per-vertex local temperature.
\end{enumerate}

\subsection{Retained local neighborhoods}

The local neighborhood structure is created by BFS. For every active root vertex
$i$ and each MISF level $\ell$, the code stores up to
\[
\texttt{nbr}[\ell] \le \texttt{num\_nbrs}
\]
other vertices together with their BFS graph distances from $i$. The arrays
\texttt{nbrs[root][level]} therefore contain truncated, local graph-distance
information, not the full all-pairs graph metric.

This point is essential: the core KK-style and FR-style updates in GRIP are
based on these truncated local neighborhoods.

\subsection{Local temperature and step normalization}

Another major difference from textbook KK and FR is that the raw force-like
vector in GRIP mostly determines \emph{direction}, not the final step size.

Let $\vect{\Delta}_i^{(t)}$ be the raw displacement computed by one of the local
kernels at round $t$. The code normalizes it to a fixed reference length
\[
e = \texttt{edge} = 32,
\]
so that the normalized direction is effectively
\[
\widehat{\vect{\Delta}}_i^{(t)} = e \, \frac{\vect{\Delta}_i^{(t)}}{\norm{\vect{\Delta}_i^{(t)}}}.
\]

Then a local heat variable $T_i$ is updated using the angle between consecutive
directions. If
\[
c_i^{(t)} =
\frac{
\widehat{\vect{\Delta}}_i^{(t)} \cdot \widehat{\vect{\Delta}}_i^{(t-1)}
}{
\norm{\widehat{\vect{\Delta}}_i^{(t)}} \,
\norm{\widehat{\vect{\Delta}}_i^{(t-1)}}
},
\]
then the code updates heat approximately as
\[
T_i \leftarrow
\begin{cases}
T_i + T_i s r c_i^{(t)}, & \text{if the direction is consistent with the previous one},\\
T_i + T_i r c_i^{(t)}, & \text{otherwise}.
\end{cases}
\]

Finally, the actual position update is effectively
\[
\vect{x}_i \leftarrow \vect{x}_i + T_i \, \frac{\vect{\Delta}_i^{(t)}}{\norm{\vect{\Delta}_i^{(t)}}},
\]
possibly multiplied by an extra factor on late FR rounds when
\texttt{final\_move\_scale\_after\_first} is active.

So in GRIP:
\begin{itemize}[leftmargin=1.5em]
\item the kernel formulas determine direction and relative balance,
\item the local heat schedule determines how far the vertex actually moves.
\end{itemize}

\section{KK-style implementations in \texttt{grip}}

\subsection{\texttt{KK\_spring\_local()}: warm-start KK for new vertices}

This routine is used only while introducing newly activated vertices at finer
levels. Each new vertex is initialized from three nearby already-placed centers,
then refined by three short local KK-style iterations.

If the three reference vertices are $j \in C_i$ with stored graph distances
$d_G(i,j)$, then the code computes
\[
\vect{\Delta}_i^{\text{KK-local}}
=
\sum_{j \in C_i}
(\vect{x}_j - \vect{x}_i)
\left(
\frac{\norm{\vect{x}_j - \vect{x}_i}^2}{d_G(i,j)^2 e^2}
- 1
\right).
\]

This is ``KK-like'' in the sense that it tries to match the geometric distance
to a target proportional to graph distance, but it is \emph{not} the classical
KK gradient.

\subsection{\texttt{KK\_spring\_v4()}: coarse/intermediate local KK}

This is the active KK-style routine on all non-finest MISF levels. If
$\mathcal{N}_{\ell}(i)$ is the retained BFS neighborhood of vertex $i$ at level
$\ell$, then the raw displacement is
\[
\vect{\Delta}_{i,\ell}^{\text{KK}}
=
\sum_{j \in \mathcal{N}_{\ell}(i)}
(\vect{x}_j - \vect{x}_i)
\left(
\frac{\norm{\vect{x}_j - \vect{x}_i}^2}{d_G(i,j)^2 e^2}
- 1
\right)
+
\vect{R}_{i,\ell}^{\text{coarse}},
\]
where the coarse repulsion term $\vect{R}_{i,\ell}^{\text{coarse}}$ is present
only in the global-repulsion engine and only when the repulsion factor is
positive.

\subsection{Exact and sampled active-set repulsion}

The additional long-range term is repulsion only, over the \emph{currently
active} vertex set $A_{\ell}$.

Define
\[
\alpha_c = \texttt{coarseFedge2}
= 0.05 \cdot \texttt{coarse\_repulsion\_factor} \cdot e^2.
\]

When the active-set size is small enough, GRIP uses exact active-set repulsion:
\[
\vect{R}_{i,\ell}^{\text{exact}}
=
\sum_{\substack{j \in A_{\ell}\\j \ne i}}
\alpha_c
\frac{\vect{x}_i - \vect{x}_j}{\norm{\vect{x}_i - \vect{x}_j}^2}.
\]

When the active set is larger, GRIP samples a subset $S_i \subset A_{\ell}\setminus\{i\}$
uniformly without replacement, with
\[
|S_i| = \min(\texttt{coarse\_repulsion\_sample}, |A_{\ell}| - 1),
\]
and rescales the result by
\[
\frac{|A_{\ell}| - 1}{|S_i|}.
\]
The sampled repulsion is therefore
\[
\vect{R}_{i,\ell}^{\text{sampled}}
=
\frac{|A_{\ell}| - 1}{|S_i|}
\sum_{j \in S_i}
\alpha_c
\frac{\vect{x}_i - \vect{x}_j}{\norm{\vect{x}_i - \vect{x}_j}^2}.
\]

This is the precise meaning of the package's ``sampled long-range interactions'':
\begin{itemize}[leftmargin=1.5em]
\item they are \emph{repulsive only},
\item they are applied over the currently active set,
\item they are exact below a threshold and sampled above it,
\item they do \emph{not} create a full KK all-pairs spring system.
\end{itemize}

\subsection{\texttt{KK\_spring\_final()}: finest-level KK with repulsion}

When \texttt{final\_mode = "kk\_repulse"}, the finest level uses
\texttt{KK\_spring\_final()} instead of FR. Its local distance-matching term is
the same as \texttt{KK\_spring\_v4()}, but now the active set is the full current
finest-level set and the repulsion scale is
\[
\alpha_f = \texttt{fedge2}
= 0.05 \cdot \texttt{repulsion\_factor} \cdot e^2.
\]

The raw displacement becomes
\[
\vect{\Delta}_i^{\text{KK-final}}
=
\sum_{j \in \mathcal{N}_{0}(i)}
(\vect{x}_j - \vect{x}_i)
\left(
\frac{\norm{\vect{x}_j - \vect{x}_i}^2}{d_G(i,j)^2 e^2}
- 1
\right)
+
\vect{R}_i^{\text{active}}(\alpha_f),
\]
where the repulsion $\vect{R}_i^{\text{active}}(\alpha_f)$ is again exact or
sampled according to the same active-set threshold logic.

\subsection{What GRIP KK is \emph{not}}

Compared with the original Kamada--Kawai algorithm, the GRIP KK-style routines
omit several defining pieces:
\begin{enumerate}[leftmargin=1.5em]
\item \textbf{No all-pairs shortest-path matrix.} The code stores only truncated
local BFS neighborhoods, not all graph distances $g_{ij}$.

\item \textbf{No explicit KK energy.} There is no energy
\[
E(X) = \frac{1}{2} \sum_{i<j} k_{ij}\bigl(\norm{\vect{x}_i - \vect{x}_j} - \ell_{ij}\bigr)^2
\]
in the implementation.

\item \textbf{No Newton--Raphson solve.} There is no per-vertex Hessian, no
largest-gradient vertex selection, and no 2-by-2 Newton solve.

\item \textbf{No full all-pairs springs.} The active force calculation is local,
with optional repulsion added separately.

\item \textbf{No weighted graph-distance KK.} Even when edge weights are present
in the graph object, the KK-style BFS distances are hop-count distances from the
unweighted adjacency traversal. The KK-style code does not read edge weights.
\end{enumerate}

So the correct description is:
\begin{quote}
GRIP implements local KK-style distance-matching kernels inside a multiscale
engine, with optional exact-or-sampled active-set repulsion. It does not
implement textbook full KK.
\end{quote}

\section{FR-style implementation in \texttt{grip}}

\subsection{Which FR routine is actually used}

The active engine calls \texttt{FR\_spring()} at the finest level when
\texttt{final\_mode = "fr"}.

There is also a function \texttt{FR\_spring\_v2()} in the codebase, but the
current engine never calls it. It appears to be an older or alternate helper
left in the source.

\subsection{Active finest-level FR kernel}

At the finest level, GRIP computes an FR-style raw displacement from three parts:
\begin{enumerate}[leftmargin=1.5em]
\item attractive force along actual graph edges,
\item repulsive force over the retained local neighborhood only,
\item optional anchor force pulling back toward the pre-final full-graph layout.
\end{enumerate}

\subsection{Attractive term}

If $(i,j) \in E$, define
\[
\ell_{ij} =
\begin{cases}
e, & \text{if the graph is unweighted},\\
e \, w_{ij}, & \text{if the graph stores edge weight } w_{ij}.
\end{cases}
\]

Then GRIP's attractive contribution from edge $(i,j)$ is
\[
\vect{A}_{ij}^{\text{GRIP-FR}}
=
(\vect{x}_j - \vect{x}_i)
\frac{\norm{\vect{x}_j - \vect{x}_i}^2}{\ell_{ij}^2}.
\]

Summing over adjacent vertices gives
\[
\vect{A}_i^{\text{GRIP-FR}}
=
\sum_{j:(i,j)\in E}
(\vect{x}_j - \vect{x}_i)
\frac{\norm{\vect{x}_j - \vect{x}_i}^2}{\ell_{ij}^2}.
\]

This term \emph{does} use edge weights, but only here. Large weights increase
the effective desired length $\ell_{ij}$ and therefore weaken the attraction.

\subsection{Repulsive term}

Let
\[
\alpha_f = \texttt{fedge2}
= 0.05 \cdot \texttt{repulsion\_factor} \cdot e^2.
\]

If $\mathcal{N}_{0}(i)$ is the retained finest-level local neighborhood of $i$,
then the FR-style repulsive term used by GRIP is
\[
\vect{R}_i^{\text{GRIP-FR}}
=
\sum_{j \in \mathcal{N}_{0}(i)}
\alpha_f
\frac{\vect{x}_i - \vect{x}_j}{\norm{\vect{x}_i - \vect{x}_j}^2}.
\]

This is the key place where GRIP differs from textbook FR:
\begin{itemize}[leftmargin=1.5em]
\item the repulsion is not over all other vertices,
\item it is only over the retained local BFS neighborhood,
\item there is no Barnes--Hut style global approximation here,
\item and there is no exact full all-pairs repulsion in the active FR stage.
\end{itemize}

\subsection{Optional anchor term}

If \texttt{final\_anchor\_factor} is positive, GRIP stores the vertex positions
at the start of the finest full-graph level and adds
\[
\vect{H}_i
=
\frac{\lambda}{e}
\left(
\vect{x}_i^{\text{anchor}} - \vect{x}_i
\right),
\qquad
\lambda = \texttt{final\_anchor\_factor}.
\]

This term is not part of original FR. It is a GRIP-specific stabilizer for the
final FR stage.

\subsection{Complete GRIP FR raw displacement}

Putting the pieces together, the active finest-level FR displacement is
\[
\vect{\Delta}_i^{\text{GRIP-FR}}
=
\vect{A}_i^{\text{GRIP-FR}}
+
\vect{R}_i^{\text{GRIP-FR}}
+
\vect{H}_i.
\]

After this, GRIP again normalizes the direction, applies its local heat-based
step-size rule, and optionally damps later rounds by
\texttt{final\_move\_scale\_after\_first} when that parameter is below $1$.

\subsection{What GRIP FR is \emph{not}}

Compared with textbook Fruchterman--Reingold, the GRIP FR stage differs in
several important ways:
\begin{enumerate}[leftmargin=1.5em]
\item \textbf{Repulsion is local, not all-pairs.} Classical FR uses repulsion
between every pair of vertices; GRIP uses only the retained BFS neighborhood.

\item \textbf{The attraction formula is not the textbook FR formula.} Original FR
uses
\[
f_a(r) = \frac{r^2}{k},
\qquad
f_r(r) = \frac{k^2}{r}.
\]
In vector form, the classical attractive vector has magnitude proportional to
$r^2/k$. In GRIP, the attractive vector has magnitude proportional to
$r^3/\ell_{ij}^2$.

\item \textbf{The step size is not controlled by a global temperature schedule.}
GRIP uses per-vertex local heat values and direction consistency instead.

\item \textbf{FR is only the finest-level kernel.} On coarser levels, the engine
uses KK-style updates, not FR.
\end{enumerate}

So the correct description is:
\begin{quote}
GRIP implements a finest-level FR-style local edge-attraction/local-repulsion
kernel inside the multiscale GRIP engine, not textbook full FR.
\end{quote}

\section{Legacy versus current default behavior}

\subsection{\texttt{grip.layout.legacy()}}

The legacy wrapper uses \texttt{grip\_layout\_adj\_cpp()} and constructs
\texttt{DrawGraph} without the extra active-set repulsion parameters. In practice
this means:
\begin{itemize}[leftmargin=1.5em]
\item coarse/intermediate levels use \texttt{KK\_spring\_v4()} with no added
active-set repulsion,
\item the finest level uses \texttt{FR\_spring()},
\item there is no coarse global repulsion term,
\item there is no final-stage anchor term,
\item there is no \texttt{kk\_repulse} final mode exposed.
\end{itemize}

\subsection{\texttt{grip.layout()} / \texttt{grip.layout.globalrep()}}

The current default wrapper uses \texttt{grip\_layout\_globalrep\_adj\_cpp()}.
Compared with legacy behavior, it adds:
\begin{itemize}[leftmargin=1.5em]
\item exact-or-sampled active-set repulsion on coarse/intermediate KK levels,
\item an optional final-stage switch between FR and KK-plus-repulsion,
\item optional final anchor and post-first-round damping for the FR final stage.
\end{itemize}

Therefore the user intuition
\begin{quote}
``there is one local KK version and another KK version with sampled long-range
interactions, but no full all-pairs KK''
\end{quote}
is essentially correct, but the precise refinement is:
\begin{itemize}[leftmargin=1.5em]
\item there are \emph{multiple} KK-style routines,
\item the extra long-range piece is repulsion only,
\item the long-range repulsion can be \emph{exact} or \emph{sampled},
\item and the package still does not implement full textbook KK.
\end{itemize}

\section{Comparison with the original algorithms}

\subsection{Original Kamada--Kawai versus GRIP KK}

\begin{center}
\small
\begin{tabular}{>{\raggedright\arraybackslash}p{0.26\textwidth} >{\raggedright\arraybackslash}p{0.31\textwidth} >{\raggedright\arraybackslash}p{0.31\textwidth}}
\toprule
\textbf{Aspect} & \textbf{Original KK} & \textbf{GRIP KK-style routines} \\
\midrule
Pair set & All vertex pairs & Truncated BFS neighborhood only \\
Graph distances & Full all-pairs shortest paths & Stored local BFS hop distances only \\
Objective & Explicit spring energy & No explicit global energy in code \\
Optimizer & Newton--Raphson per selected vertex & Direction computation plus local heat step \\
Long-range term & Built into all-pairs springs & Optional extra active-set repulsion only \\
Weights & Naturally compatible with weighted shortest paths & KK-style code does not read edge weights \\
Use in engine & Standalone layout algorithm & Coarse/intermediate kernel; optional finest-level mode \\
\bottomrule
\end{tabular}
\end{center}

\subsection{Original Fruchterman--Reingold versus GRIP FR}

\begin{center}
\small
\begin{tabular}{>{\raggedright\arraybackslash}p{0.26\textwidth} >{\raggedright\arraybackslash}p{0.31\textwidth} >{\raggedright\arraybackslash}p{0.31\textwidth}}
\toprule
\textbf{Aspect} & \textbf{Original FR} & \textbf{GRIP FR-style routine} \\
\midrule
Repulsion & All-pairs & Retained local BFS neighborhood only \\
Attraction & Edge-based, textbook FR law & Edge-based, but with GRIP's own vector formula \\
Temperature & Global cooling schedule & Per-vertex local heat adaptation \\
Weights & Usually added by variant design & Edge weights affect attraction length scale only \\
Role in engine & Standalone layout algorithm & Finest-level kernel only \\
Extra terms & None in original core & Optional anchor and late-round damping \\
\bottomrule
\end{tabular}
\end{center}

\section{Direct answer to the original question}

\textbf{Yes, essentially correct.}

The current \texttt{grip} package does \emph{not} contain a full all-pairs
Kamada--Kawai layout implementation. What it contains is:
\begin{enumerate}[leftmargin=1.5em]
\item a local KK-style warm-start routine for newly introduced vertices,
\item a local KK-style multiscale refinement routine,
\item an optional finest-level KK-style routine with exact-or-sampled active-set repulsion,
\item a finest-level FR-style routine with local, not global, repulsion.
\end{enumerate}

The ``sampled long-range interactions'' are active-set repulsive corrections,
not a full all-pairs KK spring model. Likewise, the FR stage is FR-style but not
textbook all-pairs FR.

\section{Code basis for this document}

The statements above were derived from the current source tree, especially:
\begin{itemize}[leftmargin=1.5em]
\item \texttt{R/grip\_layout.R}: public API and parameter semantics.
\item \texttt{src/grip\_rcpp.cpp}: mapping from R wrappers to \texttt{DrawGraph}.
\item \texttt{src/MishEngine.cpp}: level schedule and kernel selection logic.
\item \texttt{src/MishSupport.cpp}: KK-style and FR-style formulas, plus exact
and sampled active-set repulsion.
\item \texttt{src/DrawGraph.cpp} and \texttt{src/DrawGraph.h}: constructor
constants, neighborhood-size caps, local temperature schedule, and helper state.
\end{itemize}

\begin{thebibliography}{9}

\bibitem{kamada1989}
T. Kamada and S. Kawai.
\newblock An algorithm for drawing general undirected graphs.
\newblock \emph{Information Processing Letters}, 31(1):7--15, 1989.

\bibitem{fr1991}
T. M. J. Fruchterman and E. M. Reingold.
\newblock Graph drawing by force-directed placement.
\newblock \emph{Software: Practice and Experience}, 21(11):1129--1164, 1991.

\bibitem{gajer2002}
P. Gajer and S. G. Kobourov.
\newblock GRIP: Graph dRawing with Intelligent Placement.
\newblock \emph{Journal of Graph Algorithms and Applications}, 6(3):203--224, 2002.

\bibitem{gajer2004}
P. Gajer, M. T. Goodrich, and S. G. Kobourov.
\newblock A multi-dimensional approach to force-directed layouts of large graphs.
\newblock \emph{Computational Geometry}, 29(1):3--18, 2004.

\end{thebibliography}

\end{document}
